% SPDX-FileCopyrightText: 2025 IObundle
%
% SPDX-License-Identifier: MIT

Every unit Versat instantiates, built-in or custom, communicates with the
rest of the generated accelerator through a fixed set of interfaces.
A custom unit only needs to implement the ones it actually uses; Versat
recognizes each interface by the naming pattern of its wires, as summarized
below.

\subsubsection*{Standard Wires}

\begin{lstlisting}
input clk,
input rst,
input run,      // Asserted for a single cycle to start a run.
input running,  // Asserted for the whole duration of a run.
\end{lstlisting}

A unit may turn off internal logic while \texttt{running} is deasserted, to
save energy and speed up simulation.

\subsubsection*{Data Ports}

Inputs and outputs carry the unit's dataflow streams, one value per valid
cycle:

\begin{lstlisting}
input      [DATA_W-1:0] in0,
(* versat_latency = T *)
output     [DATA_W-1:0] out0,
input      [DELAY_W-1:0] delay0, // How many cycles to wait before in0 is valid.
output                   done,   // Only for units that take a variable amount
                                  // of time to finish; see below.
\end{lstlisting}

Since Versat must statically calculate how long it takes for data to
traverse the dataflow graph (Section~\ref{sec:data_validity}), every output
must be tagged with a \texttt{versat\_latency} attribute giving the number
of cycles between valid input and the corresponding valid output.
\texttt{DATA\_W} is the coarse-grained accelerator's data granularity;
units are not required to support more than one value of \texttt{DATA\_W},
but doing so lets Versat reuse them at different granularities.

Units whose latency cannot be known statically (most notably, units using
a Databus interface) must instead expose a \texttt{done} wire that the
accelerator polls to know when the unit has finished.

\subsubsection*{Config and State Interfaces}
\label{sec:if_config_state}

Any input wire whose name does not match one of the other interfaces below
is treated as part of the unit's \emph{Config} interface; any such output
wire is part of its \emph{State} interface:

\begin{lstlisting}
input  [ANY_SIZE:0] anyOtherInputName,  // Config
output [ANY_SIZE:0] anyOtherOutputName, // State
\end{lstlisting}

All units' configurations live together in a single register inside the
accelerator, which Versat exposes to software as a memory-mapped,
per-accelerator \texttt{Config} struct (and, symmetrically, a
\texttt{State} struct for reading unit outputs); writing a struct member
writes directly to the accelerator's memory-mapped register connected to
that wire, with no function call in between. Reading from a Config wire or
writing to a State wire is undefined behaviour and can deadlock the system;
avoiding this is the programmer's responsibility.

By default, this configuration register is shadowed: software can change
it safely while the accelerator is running, and the change is only picked
up by the units at the start of the \emph{next} run, never mid-run. This is
what lets software call \texttt{StartAccelerator} to launch a run without
waiting for it to finish, reconfigure the accelerator for the next run
while the current one is still executing, and only later call
\texttt{EndAccelerator} to wait for completion before reading any State,
memory-mapped, or accelerator-written data that the just-finished run
produced. \texttt{RunAccelerator} is a convenience wrapper that starts the
accelerator and blocks until it stops, for callers that do not need to
overlap configuration with execution.

\subsubsection*{Memory-Mapped Interface}

\begin{lstlisting}
input                 valid,
output                rvalid,
input  [DATA_W/8-1:0] wstrb,
input  [  ADDR_W-1:0] addr,
input  [  DATA_W-1:0] wdata,
output [  DATA_W-1:0] rdata,
\end{lstlisting}

This is a simple memory-mapped interface, following the IOb native
interface convention, generally used while the accelerator is not
running. It suits infrequent, large transfers, such as initializing a
lookup table, where a Databus interface (below) would be unnecessarily
heavy. Because PC emulation cannot access memory-mapped units the same way
it accesses Config/State, software must go through the functions declared
in \texttt{versat\_accel.h} rather than dereferencing pointers directly;
large transfers in particular should use \texttt{VersatMemoryCopy}, which
can use an internal DMA (limited to memory-to-accelerator and
accelerator-to-memory transfers) to speed them up.

\subsubsection*{Databus Interface}
\label{sec:if_databus}

\begin{lstlisting}
input                     databus_ready_0,
output                    databus_valid_0,
output [  AXI_ADDR_W-1:0] databus_addr_0,
input  [  AXI_DATA_W-1:0] databus_rdata_0,
output [  AXI_DATA_W-1:0] databus_wdata_0,
output [AXI_DATA_W/8-1:0] databus_wstrb_0,
output [       LEN_W-1:0] databus_len_0,
input                     databus_last_0,
\end{lstlisting}

A unit's Databus interface, similar to an AXI-lite master except that the
transfer length must be known upfront, gives it access to external RAM
while the accelerator is running; a unit can implement more than one, in
which case they are numbered starting at zero. Because a Databus transfer
does not complete in a statically known number of cycles, a unit that uses
one must implement the \texttt{done} wire, and cannot feed its data
straight into the dataflow graph: it needs internal memory to buffer
between the graph's fixed-latency timing and RAM's variable-latency
timing. The built-in \texttt{VRead} and \texttt{VWrite} units are the
standard way to read from and write to RAM this way; each keeps two
buffers and flips which one is being transferred and which one is being
streamed to/from the graph on every run, so that transfer and computation
overlap across runs (Section~\ref{sec:address_gen} covers how their
addresses are generated).

\subsubsection*{External Memory Interfaces}

\begin{lstlisting}
// Two-port
output [ANY_SIZE:0] ext_2p_addr_out_0,
output [ANY_SIZE:0] ext_2p_data_out_0,
output              ext_2p_write_0,
output [ANY_SIZE:0] ext_2p_addr_in_0,
input  [ANY_SIZE:0] ext_2p_data_in_0,
output              ext_2p_read_0,

// True dual-port (one port pair shown; _port_1 mirrors _port_0)
output [ANY_SIZE:0] ext_dp_addr_0_port_0,
output [ANY_SIZE:0] ext_dp_out_0_port_0,
input  [ANY_SIZE:0] ext_dp_in_0_port_0,
output              ext_dp_enable_0_port_0,
output              ext_dp_write_0_port_0,
\end{lstlisting}

These interfaces let a unit ask Versat to instantiate and wire up a memory
it needs, rather than declaring the memory itself: a two-port memory for
\texttt{ext\_2p}, a true dual-port memory for \texttt{ext\_dp}. Multiple
instances of either interface are numbered starting at zero, except that
for \texttt{ext\_dp} the memory index comes first and the port index
(\texttt{\_port\_0}/\texttt{\_port\_1}) comes last. The address and data
widths may differ between a memory's two ports, but the total storage they
address must match; Versat instantiates the underlying memory accordingly.

\subsubsection*{Custom Units}

Versat currently only accepts custom units written directly in Verilog.
Every interface above is optional and independent: a custom unit
implements whichever subset it needs, declares any of the parameters above
(\texttt{DATA\_W}, \texttt{AXI\_ADDR\_W}, \texttt{AXI\_DATA\_W},
\texttt{DELAY\_W}, \texttt{LEN\_W}) that its chosen interfaces require, and
can additionally declare its own parameters, which the specification
language can then set per instance.
